% ============================================================================
\chapter{同余子群上的模形式}
\label{ch:modular-forms-congruence}
% ============================================================================

% ----------------------------------------------------------------------------
\section{同余子群上的模形式与尖点形式}
% ----------------------------------------------------------------------------

\begin{definition}
设 $\Gamma$ 为 $\SL_2(\Z)$ 的一个同余子群。权 $k$ 的\textbf{模形式}
$f \in \Mk_k(\Gamma)$ 是全纯函数 $f\colon \HH \to \C$，满足：
\begin{enumerate}
  \item 对任意 $\gamma = \begin{pmatrix} a & b \\ c & d \end{pmatrix}
    \in \Gamma$，有 $f(\gamma \cdot \tau) = (c\tau + d)^k f(\tau)$；
  \item $f$ 在 $\Gamma$ 的每个尖点处全纯。
\end{enumerate}
\end{definition}

% ----------------------------------------------------------------------------
\section{$\Mk_k(\GammaO(N))$ 与 $\Mk_k(\GammaI(N))$ 的维数公式}
% ----------------------------------------------------------------------------

% TODO: 一般同余子群的维数公式（Riemann-Roch 方法）

% ----------------------------------------------------------------------------
\section{Diamond 算子与 $\langle d \rangle$ 作用}
% ----------------------------------------------------------------------------

% TODO: Diamond 算子的定义与性质

% ----------------------------------------------------------------------------
\section{oldforms 与 newforms（预告）}
% ----------------------------------------------------------------------------

% TODO: 简要预告 Atkin-Lehner 理论

% ----------------------------------------------------------------------------
\section{习题}
% ----------------------------------------------------------------------------
